\section{Conclusiones}
\label{sec:conclusiones}

El artículo \cite{piperno2025} muestra que la destilación de conocimiento con Sentence Transformers puede transferir conocimiento biomédico desde el inglés hacia el español, el francés y el alemán cuando hay pocas etiquetas. La ganancia media cae de 7,95 puntos de F1 con el 10\,\% de los datos a 0,45 con el 100\,\%, y la dirección del efecto se repite en las tres lenguas y en los tres profesores. El máximo de 12,89 puntos corresponde a una sola celda y no debe leerse como el resultado típico.

Sin embargo, la contribución es real en eficiencia de datos, no en equidad ni en validez clínica. El artículo no mide sesgo, no valida la calidad de las traducciones y no evalúa texto clínico. Su corpus son manuales educativos, sus lenguas son de altos recursos y su prueba está filtrada por similitud con el entrenamiento traducido. Esas condiciones acotan lo que puede concluirse: el método funciona en clasificación temática sobre prosa divulgativa, con un profesor inglés y un estudiante multilingüe, en el régimen de pocos datos.

Las cinco preguntas reciben respuestas convergentes. Añadir traducción y destilación se justifica con hasta el 25\,\% de las etiquetas, un profesor con conocimiento de dominio o de tarea, terminología verificada y una ventaja que supere a translate-train y a la variación entre semillas. Transferir el resultado a Colombia exige validación externa, taxonomía local, anotación clínica, baselines BETO, MarIA y TF-IDF, calibración, equidad y cumplimiento normativo. La destilación puede transferir y amplificar sesgos; reducirlos exige medirlos, intervenir datos y representaciones y auditar. La equivalencia semántica no se garantiza con MNRL ni con BLEU: se audita con léxico especializado, ida y vuelta, métricas semánticas y anotación experta.

El balance final es favorable al artículo como método y cauteloso como evidencia. Aporta un diseño replicable y una pregunta pertinente; deja abiertos la varianza, la definición de la métrica, la calidad de la traducción, el sesgo y la validación local. La destilación multilingüe merece probarse en Colombia, pero solo después de construir el corpus, la taxonomía y la evaluación que permitan saber si funciona.
